
\def\PUDcodigo{EME113}

\def\PUDcodacad{04506.24}

\def\PUDdisciplina{Processos de Fabricação I}

\def\PUDsemestre{S5}

\def\PUDhoras{80}

%NOVOS ---------------------------------------------------
\def\PUDhorasTeoria{60}
\def\PUDhorasPratica{20}

\def\PUDinterECA{---}

\def\PUDatuaisECA{---}

\def\PUDrecursos{Projetor multimídia; Quadro branco e pincel; Aulas em laboratório.}

%----------------------------------------------------------

\def\PUDcreditos{4}

\def\PUDprerequisito{ \hyperref[PUD:EME108]{EME108} }

\def\PUDpreacad{ \hyperref[PUD:EME108]{04506.20} }

\def\PUDobjetivos{Compreender a importância dos diversos tipos de processos de fabricação para a indústria;  Conhecer os tipos de processos de fabricação e suas particularidades;  Ser capaz de selecionar processos de fabricação para aplicações industriais;  Saber avaliar os produtos obtidos através dos processos de fabricação.}

\def\PUDementa{Fundição;  Laminação;  Extrusão;  Trefilação;  Forjamento;  Estampagem;  Corte e Dobra de Chapas;  Metalurgia do Pó;  Moldagem por Injeção.}

\def\PUDconteudo{ & Solidificação no processos de fundição.  Etapas do processo de fundição.  Fusão do metal.  Características metalúrgicas.  Moldes.  Desmoldagem.  Limpeza.  Rebarbação.  Controle de qualidade. 
\\ 
 & Laminação dos metais.  Esforços na laminação.  Tipos de laminadores.  Tipos de laminação.  Operações de laminação.  Produtos obtidos pelo processo de laminação. 
\\ 
 & Processo de extrusão.  Formas de extrusão.  Esforços no processo de extrusão.  Produtos obtidos pelo processo de extrusão. 
\\ 
 & Processo de trefilação.  Formas de trefilação.  Esforços no processo de trefilação.  Produtos obtidos pelo processo de trefilação. 
\\ 
 & Processo de extrusão.  Formas de extrusão.  Esforços no processo de extrusão.  Produtos obtidos pelo processo de extrusão. 
\\ 
 & Processo de estagem.  Estampagem profunda.  Matrizes para estampagem.  Esforços no processo de estampagem.  Produtos obtidos pelo processo de estampagem. 
\\ 
 & Princípios de corte e dobra de chapas.  Ferramentas de corte e dobra.  Esforços no processo de corte e dobramento de chapas.  Produtos obtidos pelo processo de corte e dobramento de chapas. 
\\ 
 & Metalurgia do pó.  Matérias-primas utilizadas na metalurgia do pó.  Mistura do pó metálico.  Sinterização.  Compactação.  Tratamentos posteriores à sinterização.  Produtos obtidos pela metalurgia do pó. 
\\ 
 & Moldagem por injeção de termoplásticos.  Termoplásticos.  Moldes para injeção de termoplásticos.  Equipamentos para injeção de termoplásticos.  Extração de peças injetadas.  Produtos obtidos pelo processo de injeção de termoplásticos. }

\def\PUDmetodo{Aulas expositórias que podem ser teóricas e/ou práticas, onde as práticas no laboratório serão marcadas no decorrer da disciplina.}

\def\PUDavalia{A avaliação será feita com aplicação de provas teóricas e/ou práticas, além da possibilidade de inclusão de trabalhos e seminários no decorrer da disciplina.}

\def\PUDbibbasica{ & \href{www.biblioteca.ifce.edu.br/index.asp?codigo_sophia=24061}{Chiaverini}, V.  Tecnologia Mecânica - Processos de Fabricação e Tratamentos.  v. 2.  2ª ed.  São Paulo, SP: Pearson Makron Books, 1986.  315p.  ISBN: 9780074500903. 
\\ 
 & \href{www.biblioteca.ifce.edu.br/index.asp?codigo_sophia=24419}{Cetlin}, P. R.;  Helman, H.  Fundamentos da Conformação Mecânica dos Metais.  2ª ed.  São Paulo, SP: Artliber, 2005.  260p.  ISBN: 8588098288. 
\\ 
 & \href{www.biblioteca.ifce.edu.br/index.asp?codigo_sophia=24424}{Silva}, A. L. C.;  Mei, P. R.  Aços e Ligas Especiais.  3ª ed.  São Paulo, SP: Blucher, 2010.  646p.  ISBN: 9788521205180. }

\def\PUDbibcomplementar{ & \href{www.biblioteca.ifce.edu.br/index.asp?codigo_sophia=39995}{Chiaverini}, Vicente.  Aços e ferros fundidos: características gerais, tratamentos térmicos e principais tipos.  7ª ed.  São Paulo, SP: ABM, 2012.  599p.  ISBN: 8586778486.
%&  BRESCIANI FILHO, E.;  Conformação Plástica dos Metais.  5ª edição.  São Paulo.  Editora UNICAMP.  1997.  386p.  ISBN 9788526801882. 
\\ 
 &  \href{www.biblioteca.ifce.edu.br/index.asp?codigo_sophia=64997}{Baldam}, R. L.; Vieira, E. A.  Fundição - Processos e Tecnologias Correlatas.  2ª ed.  São Paulo, SP: Érica, 2014.  380p.  ISBN: 9788536504469. 
\\ 
 &  \href{www.biblioteca.ifce.edu.br/index.asp?codigo_sophia=64114}{Hosford}, W. F.; Caddell, R. M.  Metal Forming – Mechanics and Metallurgy.  4ª ed.  New York, USA: Cambridge University Press, 2011.  331p.  ISBN: 9781107004528. 
\\
 &  \href{www.biblioteca.ifce.edu.br/index.asp?codigo_sophia=63248}{Groover}, Mikell P.  Introdução aos processos de fabricação.  Rio de Janeiro, RJ: LTC, 2014.  737p.  ISBN: 9788521625193.
\\
 &  \href{www.biblioteca.ifce.edu.br/index.asp?codigo_sophia=23980}{Chiaverini}, Vicente.  Tecnologia mecânica: estrutura e propriedades das ligas metálicas.  v. 1.  2º ed.  São Paulo, SP: Pearson Education do Brasil, 1986.  266p.  ISBN: 0074500899.  }

\def\PUDautor{Francisco Nélio Costa Freitas}

\def\PUDdata{2013-05-22}

\def\PUDrevisao{1}

\def\PUDrevisor{Samuel Vieira Dias}

\def\PUDdatarevisao{2019-05-15}

\PUDcorpo
